\documentclass[11pt]{report}
\usepackage[margin=1in]{geometry}
\usepackage{xcolor}
\usepackage{listings}
\usepackage{longtable}
\usepackage{booktabs}
\usepackage{enumitem}
\usepackage{hyperref}
\hypersetup{colorlinks=true, linkcolor=blue!50!black, urlcolor=blue!50!black}

\definecolor{codebg}{RGB}{245,245,248}
\lstset{
  basicstyle=\ttfamily\small, backgroundcolor=\color{codebg},
  breaklines=true, frame=single, rulecolor=\color{gray!40},
  columns=fullflexible, keepspaces=true, showstringspaces=false
}
\newcommand{\code}[1]{\texttt{#1}}

\title{\textbf{The Otacon Engine}\\[4pt]\large A small, reusable, in-house 2D game engine ---
       architecture and usage manual}
\author{}
\date{}

\begin{document}
\maketitle
\tableofcontents

\chapter{What Otacon is}
Otacon is a small 2D game engine, built as a static library
(\code{src/engine/}, namespace \code{otacon}). It is deliberately
\textbf{game-agnostic}: it knows nothing about any particular game. Three games
ship in this repository --- \emph{Canabalt} (\code{src/game/canabalt/}), the
Chromium \emph{T-Rex runner} (\code{src/game/dino/}) and \emph{Flappy~Bird}
(\code{src/game/flappy/}) --- and each links the same unchanged library, talking
to it only through public interfaces. They differ in almost every way a game can:
480$\times$320 / 600$\times$150 / 288$\times$512 logical resolutions, a scrolling
follow camera versus a fixed one, a ported generator versus an original one. A
fourth game would be one more \code{IGame} and one more CMake target; nothing in
\code{src/engine/} would move. The golden rule of the layout is: \textbf{engine
code never mentions the game, and game code never mentions GLFW / OpenGL / Vulkan
/ CoreAudio.} Everything crosses that boundary through four seams: \code{IGame},
\code{IWindow}, \code{IRenderer} and \code{IAudio}.

\section{Philosophy}
\begin{itemize}[nosep]
  \item \textbf{In-house subsystems.} Math (with a fixed-point scalar), an
        $n$-dimensional vector library, a memory manager, time management, a debug
        runtime, the OpenGL function loader, and even the PNG and CAF asset
        decoders are written from scratch. The only third-party code is the window
        library (GLFW or SDL2), fetched automatically.
  \item \textbf{Pluggable backends.} The graphics backend (legacy OpenGL, modern
        OpenGL, Vulkan) and the window backend (GLFW, SDL2) are chosen by a single
        config flag and selected at compile time --- nothing above the platform
        layer changes.
  \item \textbf{Verifiable.} The renderers are held to a byte-identical contract:
        every backend draws the same deterministic frame to the same pixels (proven
        by SHA-256). See Chapter~\ref{ch:verify}.
\end{itemize}

\chapter{Building the engine}
The engine is built as part of the project's CMake build; there is no separate
``install''. A game executable links the \code{otacon} static-library target.

\begin{lstlisting}[language=bash]
cmake -S . -B build            # first configure fetches the window library
cmake --build build -j8        # builds libotacon.a + the game exes
\end{lstlisting}

\section{Compile-time configuration}
All compile-time choices live in \code{config/build.cfg}; edit a value and re-run
CMake (it reconfigures automatically when the file changes). Any key can also be
overridden per-configure, e.g.\ \code{-DGRAPHICS\_BACKEND=VULKAN}.

\begin{longtable}{lll}
\toprule
Key & Values & Meaning \\
\midrule
\endhead
\code{GRAPHICS\_BACKEND} & \code{GL\_LEGACY} / \code{GL\_MODERN} / \code{VULKAN} & Renderer pipeline \\
\code{WINDOW\_BACKEND}   & \code{GLFW} / \code{SDL2} & Window / input / context library \\
\code{SCALAR\_TYPE}      & \code{FLOAT} / \code{FIXED} & Simulation number type \\
\code{MEMORY\_TRACKING}  & \code{ON} / \code{OFF} & Allocation tracking + leak report \\
\code{DEBUG\_RUNTIME}    & \code{ON} / \code{OFF} & Runtime debug views \\
\bottomrule
\end{longtable}

\section{Source layout}
\begin{lstlisting}
src/engine/          # ---- Otacon engine (static lib 'otacon') ----
  core/            #   Timer (scheduling), Tween, Events (Signal + EventBus)
  core/math/         #   Fixed (Q16.16), Scalar (float|fixed), Vector, Rect,
                     #   Random, Noise, Ease
  core/memory/       #   Tracked heap + Arena + Pool allocators
  core/time/         #   Clock + TimeManager (variable/fixed timestep)
  core/debug/        #   DebugRuntime (runtime view switching)
  asset/             #   In-house PNG decode/encode + CAF decode, Resources, SaveData
  audio/             #   IAudio seam, software Mixer, CoreAudio / null backend
  platform/          #   IWindow seam; glfw/ and sdl2/ backends
  render/            #   IRenderer seam; gl_modern/ gl_legacy/ vulkan/; Canvas, Atlas
  scene/             #   Entity, Node, Scene, Camera, Collision, Shapes, Raycast,
                     #   StateMachine, Emitter,
                     #   TileMap, Animator, Verlet, PathFinder, Steering,
                     #   SpatialGrid
  App.cpp IGame.hpp  #   Engine runtime + the game contract

src/game/            # ---- games on the engine (canabalt, dino, flappy) ----
src/samples/         # ---- fifteen engine samples (exe 'samples') ----
tests/smoke.cpp      # ---- headless engine self-tests ----
\end{lstlisting}

\chapter{The four seams}
\section{\code{IGame} --- the game contract}
A game is one object implementing \code{IGame}. The engine's \code{App} owns the
window, renderer, audio device, timing and debug runtime, and drives the game:

\begin{lstlisting}[language=C++]
struct GameContext {           // handed to the game once, at init
    IRenderer*    renderer;    IWindow* window;  IAudio* audio;
    DebugRuntime* debug;       int logicalW, logicalH; const char* assetDir;
};
class IGame {
    virtual void init(GameContext&) = 0;
    virtual void handleInput(const InputFrame&) = 0;   // edges + held state
    virtual void update(Real dt) = 0;                  // one simulation step
    virtual void render(IRenderer&, const DebugRuntime&) = 0;
    virtual void shutdown() {}                         // free GPU/audio handles
    virtual Color clearColor() const;  virtual const char* title() const = 0;
};
\end{lstlisting}

\code{App::run()} is the main loop: poll input, step the simulation (zero or more
fixed steps, or a clamped variable step), then render. The game is handed the
renderer to draw with and the \code{DebugRuntime} so it can honour the
clean-play / overlay state.

\section{\code{IWindow} --- window, input, context}
\code{createWindow(cfg)} returns the compiled-in backend (GLFW or SDL2). The
window exposes a uniform \code{InputFrame} (level-triggered \code{held[]} and
edge-triggered \code{pressed[]} arrays keyed by an \code{Action} enum, plus
normalised mouse position and the drag button), the GL context plumbing
(\code{makeContextCurrent} / \code{swapBuffers} / \code{glProcAddress}) and the
Vulkan plumbing (\code{vulkanInstanceExtensions} / \code{createVulkanSurface}).
\textbf{Both backends are complete} and interchangeable; selecting one is a
config flag.

\section{\code{IRenderer} --- drawing}
See Chapter~\ref{ch:draw}. A backend implements only triangle submission, the
texture entry points and the frame lifecycle; all shapes are tessellated once in
the shared \code{IRenderer} base, so every backend draws identical geometry.
Anything beyond that is an \emph{optional capability} a backend may decline ---
see \code{supportsShaders()} in Chapter~\ref{ch:samples}.

\section{\code{IAudio} --- sound}
\code{createAudio()} returns the platform audio device (CoreAudio on macOS, a
silent stub elsewhere); it never returns null, so game code is unchanged when no
device is available. The game decodes its sounds in-house and uploads PCM with
\code{createSound}, then fires one-shots with \code{play} and background music
with \code{playMusic}. See Chapter~\ref{ch:audio}.

\chapter{Writing a game}
A minimal game is an \code{IGame} that draws something each frame:

\begin{lstlisting}[language=C++]
class MyGame : public otacon::IGame {
    otacon::IRenderer* r_ = nullptr;
    void init(otacon::GameContext& ctx) override { r_ = ctx.renderer; }
    void handleInput(const otacon::InputFrame& in) override {
        if (in.isPressed(otacon::Action::Jump)) { /* ... */ }
    }
    void update(otacon::Real dt) override { /* advance simulation */ }
    void render(otacon::IRenderer& r, const otacon::DebugRuntime&) override {
        r.fillRect(100, 100, 64, 64, otacon::Color::rgb(0xf0c020));
    }
    const char* title() const override { return "My Game"; }
};
\end{lstlisting}

\code{main()} creates the \code{App}, the game, and runs:
\begin{lstlisting}[language=C++]
otacon::App app;  MyGame game;
if (app.init(cfg, &game, assetDir)) app.run();
app.shutdown();
\end{lstlisting}

\chapter{The drawing API}\label{ch:draw}
The renderer works in a fixed \textbf{logical} coordinate space ---
$480\times320$ by default, origin top-left, $y$ downward --- and each backend
applies its own projection to clip space. Colours are straight-alpha RGBA. The
public primitives:

\begin{longtable}{ll}
\toprule
Call & Draws \\
\midrule
\endhead
\code{fillRect(x,y,w,h,color)} & a filled axis-aligned rectangle \\
\code{fillRotatedRect(cx,cy,w,h,deg,color)} & a rectangle rotated about its centre \\
\code{drawRectOutline(x,y,w,h,color,thick)} & a rectangle outline \\
\code{drawLine(x0,y0,x1,y1,color,thick)} & a thick line segment \\
\code{drawImage(tex,x,y,w,h,u0,v0,u1,v1,tint)} & a texture sub-rect into a dest rect \\
\code{drawImageRotated(tex,cx,cy,w,h,deg,uvs,tint)} & a rotated textured quad \\
\code{drawText(text,x,y,scale,color)} & a string in the built-in $3\times5$ bitmap font \\
\bottomrule
\end{longtable}

\noindent Textures are created from a decoded \code{Image} with
\code{createTexture(img, repeat)} (nearest filtering, for crisp pixel art) and
freed with \code{destroyTexture}. \code{updateTexture(tex, w, h, rgba)} replaces
an existing texture's pixels at the same extent: it is the upload path for
anything rasterised on the CPU each frame, and re-using the texture object
avoids churning a GPU allocation every frame.

\textbf{The cross-backend contract:} a backend implements only
\code{submitTriangles} (solid) and \code{submitTextured}, the texture entry
points, and \code{beginFrame}/\code{endFrame}. Every shape above --- rectangles,
oriented lines, outlines, the bitmap font --- is tessellated to triangles once,
in the shared base class, so all backends rasterise byte-for-byte identical
geometry.

\chapter{The scene system}
The engine offers a small scene graph so rendering lives in one place ---
``the Scene renders the tree, nothing draws on the side.''

\begin{itemize}[nosep]
  \item \textbf{\code{Entity}} --- simulation / collision data: position, size,
        velocity, acceleration, drag, the \code{solid}/\code{fixed}/\code{moves}
        flags, a \code{scrollFactor} for parallax, and an optional texture/colour
        for the entity draw layer. \code{updateMotion} integrates it.
  \item \textbf{\code{Node}} --- a scene-tree element with \code{update(dt,cam)}
        and \code{render(r,cam)}; concrete nodes are particle emitters, sprite
        drawers, facade renderers, UI. A node may also override
        \code{renderColliders} (debug) and is skipped when \code{!visible}.
  \item \textbf{\code{Scene}} --- owns the camera, an ordered entity draw layer
        (rebuilt each frame by the game), a backdrop node list (drawn
        \emph{before} the entities, for the farthest parallax), and the main node
        list (drawn after). It ticks and renders the whole tree, then the debug
        overlays.
  \item \textbf{\code{Camera}} --- a faithful flixel follow camera: centres a
        target, lerps the scroll, clamps to bounds, and projects world to screen
        through each entity's \code{scrollFactor} (\code{screenPoint}). Also hosts
        the decaying screen-quake.
  \item \textbf{\code{Collision}} --- flixel's swept-AABB separation (\code{solveX}
        then \code{solveY}) with the same direction bookkeeping and extent
        rejection that avoids tunnelling.
  \item \textbf{\code{Emitter}} --- a particle emitter (a \code{Node}) with a
        fixed pool, burst and continuous modes, gravity/drag/spin, and an optional
        sprite sheet so particles can be solid quads or textured frames.
\end{itemize}

\chapter{Core subsystems}
\section{Math and the pluggable scalar}
\code{Real} is the engine number type. With \code{SCALAR\_TYPE=FLOAT} it is an
IEEE float (and reproduces the reference game bit-for-bit). With \code{FIXED} it
becomes \code{Fixed}, an in-house \emph{Q16.16} fixed-point type --- a
performance/determinism teaching point. (Q16.16 saturates near $\pm32768$, so an
endless runner's world-$x$ eventually overflows; that is itself the lesson about
range.) Vectors (\code{Vector.cpp}) are templated on the scalar and explicitly
instantiated for 2, 3 and 4 dimensions.

\section{Memory manager}
\code{core/memory} provides three classic allocators: a \textbf{tracked heap}
(records every live block with size/tag/call-site and prints leaks at exit), a
linear \textbf{Arena} (bump-pointer, reset per frame) and a fixed-size
\textbf{Pool} (free-list, O(1), zero fragmentation). Global \code{new}/\code{delete}
are intentionally \emph{not} overridden, so third-party allocations stay separate.

\section{Time}
\code{TimeManager} turns wall-clock time into the per-frame \code{dt}. The
\emph{variable} policy reproduces flixel exactly (\code{dt = now - last}, clamped
to \code{maxElapsed = 1/20\,s}); the \emph{fixed} policy banks real time and emits
whole fixed steps for frame-rate-independent determinism. It also drives the time
scale used by the slow-mo / fast-forward debug control.

\section{Simulation, AI and procedural helpers}
Beyond \code{Entity} and the AABB solver, the scene layer carries the subsystems
a 2D game reaches for: \code{SpatialGrid} (a uniform-grid broad phase, so
collision need not be $O(n^2)$), \code{Raycast} (ray/AABB by the slab method,
ray/tilemap by a DDA walk, and the line-of-sight query they stand in for),
\code{Verlet} (position-based dynamics for
ropes, bridges and cloth), \code{PathFinder} (A* over a \code{TileMap}, with the
heuristic exposed as the knob it really is), \code{Steering} (flocking plus
seek/flee/arrive), \code{TileMap} and \code{Animator}. \code{core/math} adds
\code{Random}, \code{Noise} and \code{Ease}.

\code{Random} is worth singling out. Before it existed, fifteen separate places
in this project had each written their own linear congruential generator, two of
them inside the engine. A generator here is a value rather than a global, so two
systems can neither perturb each other's stream nor be made irreproducible by
the order they happen to run in --- which is why the camera's quake keeps a
private one, a cosmetic shake having no business shifting the stream a level
generator reads.

\section{Debug runtime}
\code{DebugRuntime} holds independently toggleable views (colliders, camera focus,
parallax bands, grid, wireframe, velocities, HUD) plus pause/step and curated
presets. Because the views are plain state read by the engine each frame, every
backend gets identical debug visualisation for free. A single ``UI hidden'' flag
gives clean-play (and is what \code{--capture} uses for deterministic frames).

\chapter{In-house asset decoders}
\section{PNG (\code{asset/Png.cpp}, \code{Inflate.cpp})}
Rather than link an image library, the engine decodes PNGs itself.
\code{Inflate.cpp} is a from-scratch DEFLATE decompressor (RFC~1951) with its own
bit reader, canonical Huffman decoder and LZ77 window, wrapped for zlib streams
(RFC~1950). \code{Png.cpp} parses the chunk list (IHDR/PLTE/tRNS/IDAT/IEND),
inflates, undoes the per-scanline filters (None/Sub/Up/Average/Paeth) and expands
grayscale / RGB / palette+tRNS / RGBA into straight RGBA8. \code{PngWrite.cpp} is
a matching encoder (valid PNG via stored/uncompressed DEFLATE blocks --- no
compressor needed) used for screenshots and the cross-backend diff.

\section{CAF audio (\code{asset/Caf.cpp})}
\code{loadCaf} parses the (big-endian) Core Audio Format container --- the
\code{desc} chunk for the format, the \code{data} chunk for the samples --- and
decodes LPCM of any width/endianness into interleaved float in $[-1,1]$ as an
\code{AudioClip}. (Sounds in the reference game ship as $22050$\,Hz mono $16$-bit
LPCM.)

\chapter{Audio}\label{ch:audio}
The \code{IAudio} device owns an in-house software \textbf{Mixer} and a thin
platform output. The mixer holds the uploaded sounds and a pool of one-shot voices
plus one looping music voice, sums them into an interleaved stereo float buffer,
applies master + per-voice gain and a soft clamp, and resamples on upload so the
real-time render path never converts rates. A short mutex guards the voice list
between the game thread (\code{play}) and the audio thread (\code{render}). The
\textbf{CoreAudio} backend drives a DefaultOutput AudioUnit whose render callback
pulls from the mixer; the \textbf{null} backend (non-Apple) accepts everything and
stays silent. This is the only OS-specific audio code --- exactly as GLFW is the
only OS-specific window code.

\chapter{Graphics backends}
\section{The cross-backend contract}
Every backend implements only triangle rasterisation plus the frame lifecycle.
The contract:
\begin{itemize}[nosep]
  \item \textbf{Coordinates:} logical $480\times320$, origin top-left, $y$ down;
        each backend applies its own orthographic projection.
  \item \textbf{Blending:} straight alpha, \code{SRC\_ALPHA / ONE\_MINUS\_SRC\_ALPHA}.
  \item \textbf{Primitive:} triangle lists only; no backend-specific shapes.
  \item \textbf{Capabilities:} anything a backend cannot do is a question it is
        asked, never an assumption. \code{supportsShaders()} is the only one so
        far, and whatever is built on it must also work without it, so no screen
        becomes reachable on some backends and not others.
\end{itemize}

\section{Per-backend notes}
\begin{itemize}[nosep]
  \item \textbf{GL Legacy (2.1).} \code{glOrtho} + immediate-mode
        \code{glBegin/glColor/glVertex}; a compatibility context, no loader needed.
        Fixed-function throughout, so there is no programmable stage at all and
        \code{supportsShaders()} is false.
  \item \textbf{GL Modern (core 3.3).} A streamed VBO/VAO and a small
        vertex/fragment shader pair (the vertex shader does the logical$\to$clip
        projection); needs a core, forward-compatible context and a function
        loader --- shipped in-house in \code{gl\_modern/GLLoader} (no GLAD/GLEW).
        The only backend with a reachable programmable stage, so it is the one
        where \code{supportsShaders()} is true and user fragment effects compile.
  \item \textbf{Vulkan.} A full implementation: instance (MoltenVK portability),
        device, swapchain, render pass, solid + textured pipelines from embedded
        SPIR-V, per-frame dynamic vertex buffers, staging-buffer texture uploads,
        and per-frame/per-image synchronisation. Key difference, handled in the
        shader: Vulkan clip space already has $y$ pointing down, so it does
        \emph{not} flip $y$ as the GL shaders do. Validation-clean. Shaders are
        pre-compiled to SPIR-V and embedded at build time, so there is no runtime
        compiler and \code{supportsShaders()} is false; \code{updateTexture} is
        supported, via a staging buffer and a one-time command per upload.
\end{itemize}
\textbf{Vulkan on macOS} uses MoltenVK via the Homebrew loader; launch through
\code{tools/run.sh} (it sets \code{DYLD\_FALLBACK\_LIBRARY\_PATH} and
\code{VK\_ICD\_FILENAMES}). It takes the executable to launch --- \code{canabalt}
(the default), \code{dino}, \code{flappy} or \code{samples} --- so
\code{tools/run.sh samples} is how the sample gallery is run on this backend.

\chapter{Samples}\label{ch:samples}
Fifteen sample screens (\code{src/samples/}, executable \code{samples}) exercise
the engine on its own. They link \code{otacon} and nothing else, and generate
every texture they use in code --- a sample that borrowed a game's art would mean
a seam had leaked. One gallery object hosts them all, so adding a sample is one
\code{.cpp} plus one row in \code{Registry.cpp}.

\section{Per-pixel work and the shader seam}
Two samples (2D lighting, and the shader playground) need per-fragment control,
which sits awkwardly against the cross-backend contract: only the modern OpenGL
backend has a reachable programmable stage. The engine therefore offers both
halves and lets a sample choose:

\begin{itemize}[nosep]
  \item \code{updateTexture(tex, w, h, rgba)} re-uploads an existing texture's
        pixels. A sample shades a buffer in plain C++ and pushes it once per
        frame. This works on every backend, GL legacy included.
  \item \code{createEffect(fragmentSrc, log, logSize)} compiles a user fragment
        shader, with \code{useEffect} and \code{setEffectUniform} to drive it.
        The engine keeps the vertex stage so the projection stays identical.
        \code{supportsShaders()} reports whether the backend can do this at all:
        true for GL modern, false for GL legacy (fixed-function) and for Vulkan
        (pre-built SPIR-V, no runtime compiler).
\end{itemize}

Every sample that uses an effect also carries the CPU path, so no sample is
backend-exclusive and Chapter~\ref{ch:verify}'s contract still holds everywhere.

\chapter{Debug \& dev tools}
These are engine-level and work in any game (state shown in the HUD when visible):
\begin{itemize}[nosep]
  \item \textbf{Perf + memory overlay (\code{H}).} Frame-time graph (60/30\,fps
        lines), live draw-call/vertex counts read from the renderer, and the
        memory manager's live/peak bytes --- all backend-agnostic.
  \item \textbf{Time scale (\code{+}/\code{-}).} Scales the \emph{input} time
        $0.1\times$..$4\times$, so the fixed-timestep \code{dt} is unchanged and
        slow-mo keeps determinism.
  \item \textbf{Screenshots (\code{F12}) + capture.} The renderer reads the
        framebuffer back and writes it through the in-house PNG encoder.
        \code{--capture <png> --frames N} renders \code{N} deterministic frames
        (UI hidden) and quits --- used by \code{tools/backend-diff.sh}.
  \item \textbf{Clean play (\code{`}).} Hides all dev UI at once.
\end{itemize}

\chapter{Verification}\label{ch:verify}
The renderers are held to the byte-identical contract by \code{tools/backend-diff.sh},
which renders the same deterministic frame in each graphics backend and compares
the PNGs. With both window backends in play, the full matrix --- \{GLFW, SDL2\}
$\times$ \{GL\_MODERN, GL\_LEGACY, VULKAN\} --- produces \textbf{the same SHA-256
across all six combinations}: the rigorous proof that neither the window seam nor
the graphics seam perturbs a single pixel. The headless \code{otacon\_smoke} target
additionally unit-tests the integrator, collision, the emitter, the fixed-point
math, the PNG/CAF decoders and the audio mixer.

\end{document}
